\subsubsection{Where is the prescription upload data coming from?}
\label{subsec:upload-data-sources}
The report's upload analyses use
\texttt{EDLDB.BT\_HCA\_HCDM.rxp\_upload\_entity}.
Table~\ref{tab:upload-source-resource-types} compares the entries in that
source with \texttt{rxp\_notification\_service.upload\_entity}, resolved in
the \texttt{EDLDB} database, by \texttt{associated\_resource\_type}.

This verification covers all rows currently available in each source,
without date, CMPU, vet-diet, upload-type, or study-population filters.
In Table~\ref{tab:upload-source-resource-types}, each physical row counts
once; no deduplication or joins are applied.
These source counts therefore describe the source tables rather than the report's
study cohort. They do not establish record-level equivalence between the
two sources or document their upstream lineage.

\paragraph{How many upload-source entries have prescription IDs for each associated resource type?}
Upload-source entries and entries with a non-zero prescription ID, by associated resource type.

\begin{longtable}{@{}lrrrr@{}}
\caption{Upload-source entries and entries with a non-zero prescription ID, by associated resource type.}
\label{tab:upload-source-resource-types}\\
\toprule
 & \multicolumn{2}{c}{\shortstack{\texttt{EDLDB.BT\_HCA\_HCDM.}\\\texttt{rxp\_upload\_entity}}} & \multicolumn{2}{c}{\shortstack{\texttt{rxp\_notification\_service.}\\\texttt{upload\_entity}}} \\
\cmidrule(lr){2-3}\cmidrule(lr){4-5}
\texttt{associated\_resource\_type} & All entries & With Rx ID & All entries & With Rx ID \\
\midrule
\endfirsthead
\toprule Resource type & HCDM total & With Rx ID & Notification total & With Rx ID \\
\midrule
\endhead

{DIET\_PRESCRIPTION} & 181,943 & 181,943 & 184,637 & 184,635 \\

{DRUG\_PRESCRIPTION} & 45,458 & 45,458 & 45,551 & 45,551 \\

{[SQL NULL]} & 1 & 0 & 41,475 & 0 \\
\midrule
{Total} & 227,402 & 227,401 & 271,663 & 230,186 \\
\bottomrule
\end{longtable}


\noindent\textit{Extraction:} 2026-11-02 16:02:32 UTC.


\noindent\textit{Notes:} Rows include every resource type present in either
source. A zero means that the source has no entries for that type.
\texttt{[SQL NULL]} denotes a missing resource type, separately from any
literal string value. Totals are queried independently from each source
and reconciled against the sum of its grouped counts before publication.
The extraction time above records the snapshot used for this comparison;
rerunning the query may change the counts as sources are updated.
\textit{With Rx ID} counts entries with at least one populated, non-zero
prescription identifier in HCDM \texttt{RX\_ID} or notification-service
\texttt{ASSOCIATED\_RESOURCE\_IDS}. Braces, array delimiters and quotes are
removed before checking individual comma-separated IDs; SQL NULL, JSON null,
blank, empty-list and numeric-zero values are excluded. An entry counts once
regardless of how many IDs it contains. IDs are checked for presence, without
joining to the prescription registry. Both count columns reconcile to their
separately queried source totals.
This table is refreshed independently of the subsequent order analyses;
each analysis reports its own extraction time.


% Generate from retained aggregates with python -m shared.python.analyze_upload_source_time_coverage.
\Needspace{20\baselineskip}
Table~\ref{tab:upload-source-time-coverage} shows the minimum and maximum creation timestamps of the entries currently available in each source, by resource type and overall. All timestamps are shown in UTC.
\paragraph{What creation-date ranges do the upload sources cover?}
Creation-time coverage of current upload-source entries (UTC).

\begin{longtable}{@{}lcccc@{}}
\caption{Creation-time coverage of current upload-source entries (UTC).}
\label{tab:upload-source-time-coverage}\\
\toprule
 & \multicolumn{2}{c}{\shortstack{\texttt{EDLDB.BT\_HCA\_HCDM.}\\\texttt{rxp\_upload\_entity}}} & \multicolumn{2}{c}{\shortstack{\texttt{rxp\_notification\_service.}\\\texttt{upload\_entity}}} \\
\cmidrule(lr){2-3}\cmidrule(lr){4-5}
Resource type & Minimum created & Maximum created & Minimum created & Maximum created \\
\midrule
\endfirsthead
\toprule Resource type & HCDM min & HCDM max & Notification min & Notification max \\
\midrule
\endhead

{DIET\_PRESCRIPTION} & \shortstack{2026-07-30\\19:51:39} & \shortstack{2026-11-04\\04:36:23} & \shortstack{2026-07-29\\17:23:29} & \shortstack{2026-11-04\\13:20:17} \\

{DRUG\_PRESCRIPTION} & \shortstack{2026-10-01\\14:46:41} & \shortstack{2026-11-03\\20:49:10} & \shortstack{2026-10-01\\14:46:41} & \shortstack{2026-11-03\\20:49:10} \\

{[SQL NULL]} & \shortstack{2026-11-04\\04:45:37} & \shortstack{2026-11-04\\04:45:37} & \shortstack{2026-07-30\\19:49:25} & \shortstack{2026-11-04\\13:21:37} \\
\midrule
{Overall} & \shortstack{2026-07-30\\19:51:39} & \shortstack{2026-11-04\\04:45:37} & \shortstack{2026-07-29\\17:23:29} & \shortstack{2026-11-04\\13:21:37} \\
\bottomrule
\end{longtable}


\noindent\textit{Notes:} No date, order-ID or prescription-ID filters are applied. HCDM uses \texttt{CREATED\_TS\_UTC}; notification-service \texttt{CREATED\_AT} is converted to UTC before aggregation. Minimum and maximum ignore missing timestamps. Displayed values are rounded to the nearest second (half seconds round upward); the retained extract preserves full precision. Missing creation timestamps: 0 in HCDM and 0 in the notification-service source. The ranges describe available entries, without establishing continuous coverage or the full history of either system. Extraction: 2026-11-02 16:09:54 UTC.



\clearpage

